\chapter{Listados de los certificados en formato PEM}
Como exige el enunciado de la práctica (\emph{``se incluirán listados de todos los ficheros obtenidos en formato texto, hexadecimal, base64, PEM o similar''}), se reproducen a continuación los certificados utilizados en la práctica en su formato PEM, es decir, su representación textual base64 delimitada por las cabeceras \texttt{-----BEGIN CERTIFICATE-----} y \texttt{-----END CERTIFICATE-----}.

\section{Certificado RAÍZ de la asignatura (\texttt{certificadoRaiz.crt})}
Es el certificado autofirmado de la autoridad de certificación de la asignatura de \emph{Seguridad de la Información} (CN \texttt{Asignatura Seguridad de la Informacion}, ULPGC). Es el certificado que se ha incorporado al almacén del sistema en el capítulo correspondiente para que el cliente confíe en cualquier certificado emitido por esta CA, en particular en el certificado de servidor generado para \texttt{www.ejemplo.com}.

\begin{lstlisting}[basicstyle=\ttfamily\scriptsize, caption={certificadoRaiz.crt en formato PEM}]
-----BEGIN CERTIFICATE-----
MIIEAzCCAusCFHxge6KeuU3jHyJGgu/ZsLCaKfm2MA0GCSqGSIb3DQEBCwUAMIG9
MQswCQYDVQQGEwJFUzETMBEGA1UECAwKTGFzIFBhbG1hczEjMCEGA1UEBwwaTGFz
IFBhbG1hcyBkZSBHcmFuIENhbmFyaWExDjAMBgNVBAoMBVVMUEdDMQwwCgYDVQQL
DANFSUkxLzAtBgNVBAMMJkFzaWduYXR1cmEgU2VndXJpZGFkIGRlIGxhIEluZm9y
bWFjaW9uMSUwIwYJKoZIhvcNAQkBFhZpbmZvQG9wZW5zc2wuY2ljZWkuY29tMB4X
DTIzMDIyODA5MzU0MloXDTQzMDIyMzA5MzU0Mlowgb0xCzAJBgNVBAYTAkVTMRMw
EQYDVQQIDApMYXMgUGFsbWFzMSMwIQYDVQQHDBpMYXMgUGFsbWFzIGRlIEdyYW4g
Q2FuYXJpYTEOMAwGA1UECgwFVUxQR0MxDDAKBgNVBAsMA0VJSTEvMC0GA1UEAwwm
QXNpZ25hdHVyYSBTZWd1cmlkYWQgZGUgbGEgSW5mb3JtYWNpb24xJTAjBgkqhkiG
9w0BCQEWFmluZm9Ab3BlbnNzbC5jaWNlaS5jb20wggEiMA0GCSqGSIb3DQEBAQUA
A4IBDwAwggEKAoIBAQDgNHLpTIW6K1s/l3w5ExEhoYGtSMnKAMP9uO+++SYvmclM
PWL4whNRG/SNoAiQb6MQ+D/N8qlgPT0KSUuPiiFJfrPkyVZXZiCRTQajRWq7T5wz
EPA1lGndbE36CDh97eqZSdtRWHjOM3iWJNW/jgH6+RkLIxnZA0mlHqC58WWGYbJQ
zqIm/dIXDhgYNCVZUyd840L3vAM1afRwGImbEtwIttYBxkIOhIpXT+/dijKffK1y
mRgn1cZ3JWZAsM4JevsEd4ZNZBr29qMwQT4kjq+dYBTrGfjYhm2rqFnHP6OtUvze
t3jaULqKKJ59oUUF6eMOGkGUb6PwKMqTO+aX+h+LAgMBAAEwDQYJKoZIhvcNAQEL
BQADggEBAI8F+oAf68Gv7+e9uuPAiD5PRbkjNmCuDx4gyHynO1BkClQAY1TtLiCJ
vROf2ZXkOR0Eah33vHz+Do5cf5KnMFVtFlgQSmn1HLSSL0V+xmlaKAUy7L4S/esy
+DD0VFXKNAFVkt6cU2SUu2vaZTGT+JU99lJT2t6G5i7MVYtImkCgvZEvvCfU51j2
rZ/DEONYdmjLJfLl2wQuGzC0k7QPNmsm+f3TC+Stz1pA1kkvTzxifGpSiDAY/W9d
IYa0Vkzo1sDGalQYsuDIgcQ/eB29WX9FrC1sjt8k90uRj3n2bi69BVok8aucpERy
5K0bVyrCettoaP67iKN3JR6uoLe6Lm0=
-----END CERTIFICATE-----
\end{lstlisting}

\section{Certificado del servidor (\texttt{servidor.crt})}
Es el certificado generado para los dominios \texttt{www.ejemplo.com} y \texttt{ejemplo.com}, firmado por la CA de la asignatura tal como se describe en el capítulo de creación del certificado. Su \emph{Subject} es \texttt{/C=ES/ST=LasPalmas/L=LasPalmas/O=Ejemplo/CN=www.ejemplo.com} y contiene la extensión \texttt{X509v3 Subject Alternative Name} con \texttt{DNS:www.ejemplo.com, DNS:ejemplo.com}.

\begin{lstlisting}[basicstyle=\ttfamily\scriptsize, caption={servidor.crt en formato PEM}]
-----BEGIN CERTIFICATE-----
MIIE/TCCA+WgAwIBAgIUIXiViJhmcXJPJytlLmeJEJpkbicwDQYJKoZIhvcNAQEL
BQAwgb0xCzAJBgNVBAYTAkVTMRMwEQYDVQQIDApMYXMgUGFsbWFzMSMwIQYDVQQH
DBpMYXMgUGFsbWFzIGRlIEdyYW4gQ2FuYXJpYTEOMAwGA1UECgwFVUxQR0MxDDAK
BgNVBAsMA0VJSTEvMC0GA1UEAwwmQXNpZ25hdHVyYSBTZWd1cmlkYWQgZGUgbGEg
SW5mb3JtYWNpb24xJTAjBgkqhkiG9w0BCQEWFmluZm9Ab3BlbnNzbC5jaWNlaS5j
b20wHhcNMjYwMzEyMTczODU1WhcNMjcwMzEyMTczODU1WjBhMQswCQYDVQQGEwJF
UzESMBAGA1UECAwJTGFzUGFsbWFzMRIwEAYDVQQHDAlMYXNQYWxtYXMxEDAOBgNV
BAoMB0VqZW1wbG8xGDAWBgNVBAMMD3d3dy5lamVtcGxvLmNvbTCCASIwDQYJKoZI
hvcNAQEBBQADggEPADCCAQoCggEBALyc7f5LjOXXli8I6ylBR1pGJw2jLlJ2y9d4
MfwPXxDi3YuyOwVR9y6oNYfOurd7qf3H5UO/vuN12nBWqaeBbKhMjjfA1vpcEsaO
6R2+x/ehpwLnABSFxYZRs7lukONsQQsDia6aDlCTYHBV6eHZzsFcuQvStVeEz4AI
UGof5FORe+TnZixF9smkK+Q0Wy3DiqaX9g5tDdskLjnF3hwNWGfIaDwdqWbc66XY
zYFPKtdZNW9CZGK7+Ni6d60ZNQ3MqtrO+wFXCfWnANyGTuz2lq5Dc1PGotqI4oYS
voNLegmnbjR2fMuZpDbiqGZmwiz/GNKYNveoCqxp6nWSgct1IwUCAwEAAaOCAU4w
ggFKMCcGA1UdEQQgMB6CD3d3dy5lamVtcGxvLmNvbYILZWplbXBsby5jb20wCQYD
VR0TBAIwADALBgNVHQ8EBAMCBeAwHQYDVR0OBBYEFNKsb5pmruQ1ZAffiEZg2C4K
0kEJMIHnBgNVHSMEgd8wgdyhgcOkgcAwgb0xCzAJBgNVBAYTAkVTMRMwEQYDVQQI
DApMYXMgUGFsbWFzMSMwIQYDVQQHDBpMYXMgUGFsbWFzIGRlIEdyYW4gQ2FuYXJp
YTEOMAwGA1UECgwFVUxQR0MxDDAKBgNVBAsMA0VJSTEvMC0GA1UEAwwmQXNpZ25h
dHVyYSBTZWd1cmlkYWQgZGUgbGEgSW5mb3JtYWNpb24xJTAjBgkqhkiG9w0BCQEW
FmluZm9Ab3BlbnNzbC5jaWNlaS5jb22CFHxge6KeuU3jHyJGgu/ZsLCaKfm2MA0G
CSqGSIb3DQEBCwUAA4IBAQA/0THV7MY2APaXBtuqxougkW1aDzOGFU2TmzCZxI9b
kqmd1eHJX8sD6zEi1cQpDRMyIyrE2Dp1ztzETgBUjTP1lyPh3lrvYvYa1gN9XeMu
+g8um7qvsaHTX7CFZ0d9NKmTWzdsSi3M05OfinEsLSDe4tiXrl+fMeed2qMQS7Kp
mccU81V0NEVvm6u7IrP3Hyf56eAa/kgrdwpXUWdxubHR9+XhcPShCih6A2O35G48
+8pmsmpoIQZFD4wjwY6i6xz2UkixT8hpJFkoPcSG14MgpjssDWgxQf+xAD0CF2GS
++r8S7+IdE06gVo3dKFQGJeU02FXRv8wDSP4B0itZy+J
-----END CERTIFICATE-----
\end{lstlisting}

\section{Verificación de la cadena}
Con ambos ficheros disponibles se puede verificar de forma independiente que el certificado del servidor está correctamente firmado por la CA de la asignatura:
\begin{lstlisting}[language=bash]
openssl verify -CAfile certificadoRaiz.crt servidor.crt
# servidor.crt: OK
\end{lstlisting}
La salida \texttt{OK} confirma que el certificado del servidor pertenece a la cadena de confianza encabezada por el certificado RAÍZ de la asignatura, y que en consecuencia cualquier cliente que reconozca dicha CA aceptará sin avisos las conexiones HTTPS al servidor \texttt{www.ejemplo.com}.
